\begin{abstract}

% The increasing number of digital regulations and a reality in which compliance
% with them is difficult to always achieve, has led to an increasing interest
% in runtime enforcement solutions that can automatically ensure compliance
% with regulations.
% However, modern runtime enforcement tools require regulations to be formalized.
% The Lex programming language was designed as a highly readable and audit-able
% formalization language.
% This thesis brings the concept of refinement to Lex as a method to allow for
% maximal code reuse while allowing the flexibility to adapt specifications
% for different use cases.

The growing number of digital regulations, combined with the difficulty of consistently ensuring compliance, has led to increased interest in runtime enforcement techniques that automatically check and enforce regulatory conformance.
However, such approaches rely on formalized representations of legal norms.
Lex is a domain-specific language designed to provide a highly readable and auditable way to formalize legal specifications.
This thesis extends Lex with a notion of refinement, enabling modular specification development through code reuse and incremental adaptation.
Refinement allows abstract legal specifications to be concretized for specific use cases while maintaining a clear connection to their original intent.

\end{abstract}
